\section{Construction}\label{sec:tsps-construction}

In brief, the construction proceeds as follows. The signing key of an underlying
structure preserving signature scheme is secret-shared among $n$ signers using
the MPC protocol of \Cref{sec:tsps-mpc}. Key generation and verification are
then generic: key generation is a direct MPC evaluation of $\SPSgen$, and
verification needs no protocol at all, since $\SPSverify$ takes only public
inputs. Signing is where the construction has to work for its keep.  Naively
evaluating $\SPSsign(\sk, \vec{\msg})$ under $\algo{MPC}$ would still force the
signers to interact once $\vec{\msg}$ is known, since $\vec{\msg}$ would simply
be one more input to whatever interactive calls the evaluation makes. We instead
identify a decomposition of $\SPSsign$, satisfied by both schemes of
\Cref{sec:tsps-instantiation}, under which every interactive call can be pushed
into an offline phase run before $\vec{\msg}$ is known, leaving an online phase
that each signer completes alone. This decomposition, not the use of
$\algo{MPC}$ as such, is the technical content of this section.

We assume throughout that the signers agree in advance on the parameters $(\ell,
t, n)$, and that $\algo{MPC}$ is available to them.  \Cref{sec:tsps-evaluation}
discusses two concrete realisations.

\paragraph{Setup.} During setup, the signers fix the vector length $\ell$, the
corruption threshold $t$, and the number of signers $n > t$, together with the
description of the pairing groups over which the underlying structure preserving
signature scheme is defined.

\paragraph{Key generation and verification.} Both are generic. Neither depends
on the message, and neither requires anything beyond running the corresponding
algorithm of $\SPS$ once, under $\algo{MPC}$ in the case of key generation, and
exactly as written in the case of verification.

Key generation produces a public key distributed exactly as one produced by
$\SPSgen$, together with a $(t+1, n)$ sharing of the corresponding secret key.
Every sampling step of $\SPSgen$ becomes a call to $\algo{RandShare}$, and every
further field or group element it derives becomes a call to $\algo{Mul}$
followed by local operations. Lifting the resulting secret shares to the group
elements of $\pk$ is local by linearity preservation. Only $\pk$ is ever
reconstructed. On input $(\secparam, \ell, t, n)$, the signers do the following:

\begin{enumerate}

  \item Invoke $\algo{RandShare}$ once for each field element of $\sk$, and
$\algo{Mul}$ for each further secret value $\SPSgen$ derives from them, to
obtain $(t+1, n)$ shares $\shr{\sk}$ of the full secret key.

  \item Each signer locally lifts its shares to the group elements comprising
$\pk$, computing $\shr{a} = \gen^{\shr{\alpha}}$ in $\group{1}$ and
$\shr{\widetilde{a}} = \ggen^{\shr{\alpha}}$ in $\group{2}$ as required.

  \item Call $\algo{Open}$ in each source group to reconstruct and publish
$\pk$. Each signer stores its share of $\shr{\sk}$, together with any shares of
derived values that later phases will need.

\end{enumerate}

If any call to $\algo{MPC}$ returns $\bot$, the signers abort. Verification is
that of the underlying scheme, unchanged: on input $(\pk, \vec{\msg}, \SPSsig)$,
the verifier returns $\SPSverify(\pk, \vec{\msg}, \SPSsig)$, so that
$\TSPSverify = \SPSverify$. Since verification touches no secret share, there is
nothing here for $\algo{MPC}$ to do.

\paragraph{Why signing is not just another generic evaluation.} Write
$\SPSsign(\sk, \vec{\msg})$ as a computation on the secret key, on randomness
sampled during signing, and on the message vector. A generic secure evaluation
of this computation would supply $\vec{\msg}$ as an ordinary input, alongside
$\shr{\sk}$, to whatever sequence of $\algo{Mul}$, $\algo{Inv}$, and
$\algo{ExpGrp}$ calls $\SPSsign$ requires. Since $\vec{\msg}$ is not known until
signing time, every one of those calls would have to happen after the message
arrives, and every one of them is interactive. evaluating $\SPSsign$ generically
under $\algo{MPC}$ does not, on its own, give an online phase free of
communication. Precomputing the calls that never touch $\vec{\msg}$ is standard,
and gives some offline/online split for essentially any MPC-based signing
protocol. What is not standard is arranging for every single call that does
touch $\vec{\msg}$ to be avoidable, so that the online phase contains no
interactive call at all. That depends on the algebra of $\SPSsign$, not on
$\algo{MPC}$, and need not hold for an arbitrary signature scheme.

We therefore require, of $\SPS$, that $\SPSsign(\sk, \vec{\msg})$ decomposes
into two parts. The message-independent part consists of samplings from $\Zp$
and of field and group operations that $\algo{MPC}$ provides. The
message-dependent part uses only operations the signers perform locally, namely
products of shares and exponentiation of a public group element by a
secret-shared field element. Both schemes of \Cref{sec:tsps-instantiation}
satisfy this. The messages of a structure preserving signature scheme are
themselves group elements, public at signing time. The linearity preservation
property of \Cref{sec:tsps-mpc}, applied to the message vector, is therefore
exactly what supplies the second condition, and that is what makes the online
phase non-interactive. \Cref{sec:tsps-instantiation} checks this decomposition
explicitly for each of the two underlying signature schemes.

\paragraph{Offline signing.} The offline phase performs, ahead of time, every
part of signing that does not depend on the message. Because $\SPSsign$ samples
its randomness independently of the message, and because the values derived from
that randomness are likewise message-independent, this phase absorbs all of the
interaction that threshold signing requires. It can be run repeatedly to build
up a store of pre-processed material, one entry per future signature.

On input the secret key shares, the signers first invoke $\algo{RandShare}$ to
obtain shares of the randomness that $\SPSsign$ samples. They then invoke
$\algo{Mul}$, $\algo{Inv}$, and $\algo{ExpGrp}_1$ as the scheme requires,
interleaved with local liftings to the source groups and local products of
shares. The result is a sharing of every value the signing algorithm derives
from the secret key and that randomness.  Those components of the signature that
do not depend on the message are complete at the end of this phase, and are
stored as $\shr{\state}$ along with the shares the online phase consumes.  If
any call to $\algo{MPC}$ returns $\bot$, the signers abort.

Each stored entry is consumed by exactly one invocation of $\algo{OnSign}$ and
then discarded. Reusing an entry across two messages would reuse the signing
randomness, which is fatal for both schemes of \Cref{sec:tsps-instantiation}.
In each case the message-dependent component of the signature is a deterministic
function of the message and the pre-processed shares. Two signatures under the
same entry therefore yield a relation between the message vectors and the secret
key. The single-use requirement is reflected in the session identifiers of
Experiment~\ref{exp:tsps-adp-ts-uf-1}. In practice the store must be maintained
as a queue, with entries removed on consumption rather than cached.

\paragraph{Online signing.} Online signing takes the message vector and produces
a signature share. By the decomposition above, the only operations needed at
this point are products of shares and exponentiation of a public group element
by a secret share, both of which each signer performs locally.  No signer
communicates with any other during this phase, and its cost is comparable to
that of an ordinary invocation of $\SPSsign$.

On input a message vector $\vec{\msg} \in \group{1}^\ell$, each signer combines
its share of the pre-processed material with $\vec{\msg}$. It raises the public
group elements of $\vec{\msg}$ to its shares of the corresponding field elements
and takes the product, producing its signature share $\shr{\SPSsig}$.  Both
operations are local by linearity preservation, so this phase makes no call to
$\algo{MPC}$ at all.

\paragraph{Combining.} Combining reconstructs the signature from the shares.  It
is performed by whichever party is to hold the signature, and needs no
participation from the signers beyond their having published their shares. In
the benchmarks of \Cref{sec:tsps-evaluation} this party is one of the signers,
which is why the reported online timings include the cost of combining.

On input $t+1$ signature shares, the combining party calls $\algo{Open}$ in each
source group to reconstruct each group element of the signature from the
corresponding shares, and returns $\SPSsig$. If any call returns $\bot$, it
returns $\bot$ instead.

\paragraph{Robustness and abort.} In the honest majority setting, where $t <
n/2$, the signature can always be reconstructed from the shares of the honest
signers alone, so the scheme is robust. In the dishonest majority setting, where
$t \geq n/2$, robustness is no longer guaranteed, and in the extreme case $t =
n-1$ the sharing is additive between all $n$ signers, so every signer must
contribute for combining to succeed. If the realisation of $\algo{MPC}$ supports
identifiable abort, however, then so does the threshold signature scheme: a
corrupt signer that contributes a bogus share causes an abort in which it is
identified, and can be excluded from subsequent signing.

\begin{lemma}\label{lem:tsps-construction-correctness} Let $\SPS$ be a correct
structure preserving signature scheme whose signing algorithm admits the
decomposition described above, and let $\algo{MPC}$ satisfy the properties of
\Cref{sec:tsps-mpc}. Then the threshold structure preserving signature scheme
obtained from $\SPS$ by the construction above is correct.  \end{lemma}

\begin{theorem}\label{thm:tsps-construction-security} Let $\SPS$ be a structure
preserving signature scheme satisfying EUF-CMA security whose signing algorithm
admits the decomposition described above, and let $\algo{MPC}$ satisfy the
properties of \Cref{sec:tsps-mpc}. Then, for any corruption bound $|\Cor| \leq
t$ and any PPT adversary $\adv$ against the threshold structure preserving
signature scheme obtained from $\SPS$ by the construction above, there is a PPT
adversary $\advB$ such that
  %
  \begin{equation*}
  %
    \Adv_{\TSPS, \adv}^{\text{adp-TS-UF-1}}(\secparam) \leq \binom{n}{t} \cdot
\Adv_{\SPS, \advB}^{\text{EUF-CMA}}(\secparam),
  %
  \end{equation*}
  %
so the scheme is adaptive TS-UF-1 secure whenever $n$ is polynomial in
$\secparam$ and $t$ is constant, or whenever the loss is absorbed as described
in \Cref{sec:tsps-proof-1}.  \end{theorem}
